% abstract.tex : Hebrew abstract for IEEE paper
% Dr. Yael Mizrahi : LaTeX Technical Author

\selectlanguage{hebrew}


מאמר זה מציג מערכת ניווט ביומימטית לרחפנים זעירים בסביבות נטולות GPS, בהשראת אקולוקציית עטלפים ואינטגרציה רב-חושית במוח העטלף. המערכת משלבת מדידות טווח ודופלר מסונאר אולטראסוני בתדר 40~קילו-הרץ, נתוני \en{IMU} בקצב 200~הרץ, וזרימה אופטית ממצלמה ברזולוציה נמוכה, במסגרת מיזוג חיישנים הטרוגנית המבוססת על מסנן קלמן מורחב (\en{EKF}) מודע-דופלר. מסגרת ה-\en{SLAM} הביומימטית משתמשת ב-\en{FastSLAM 2.0} עם סיוע דופלר למיפוי עקבי בסביבות פנימיות. תוצאות סימולציית \en{Monte Carlo} (50~ריצות) מראות שיפור של 40~אחוזים בדיוק אמידת המהירות בהשוואה לגישות \en{EKF} סטנדרטיות. אימות ניסויי על רחפן \en{Crazyflie 2.1} במשקל 250~גרם מדגים שגיאת מיקום מקסימלית של 0.3~מטרים בטיסות של 5~דקות בסביבות פנימיות נטולות \en{GPS}. צריכת ההספק הכוללת של המערכת היא 95~מיליוואט, המתאימה לרחפנים זעירים. התוצאות מאמתות את היתכנות הגישה הביומימטית לניווט אוטונומי בתנאי חוסר \en{GPS}.
